\section{Perception}
\subsubsection*{Things to Write}

\begin{itemize}
  \item If you are targeting both challenges, you must include separate discussions of how your perception system matches each course.
  \item Describe processing of the sensor data to identify course features and items around the vehicle, including obstacles and lane markings, and the use of machine vision to perceive the environment.
  \item If targeting AutoNav, describe how the system identifies lanes and obstacles.
  \item If targeting SelfDrive, describe how the vehicle senses pedestrians, stop signs, and road obstacles.
  \item Describe the internal representation of the course built from the perception data which is used by the driving logic. This might be a sophisticated internal representation of the environment or a simple list of features. Include a discussion of how it is created and updated as the robot moves.
  \item Clearly identify how your design of the sensor data processing system was driven by the perception requirements you listed.
  \item For each perception requirement you listed, compare the target value with the most recentlymeasured actual value, and discuss what this tells you about your vehicle
  \item Developed a custom C++ Iterative Closest Point (ICP) node to process LiDAR point clouds for robust spatial mapping.
  \item Implemented a 3D-to-2D vertical slicing filter (using ROS 2 pointcloud iterators) to compress 3D LiDAR data, isolating course obstacles at specific Z-heights while filtering out ground and ceiling noise to prevent false obstacle detection.
\end{itemize}

\subsection{Object Detection Pipeline}
 
Object detection on the Self Drive course relies on a custom five-class
YOLO11 model that recognizes barrels, stop signs, tires, mannequins,
and potholes. The model runs on the vehicle in a ROS~2 node paired with
a ZED~2 stereo camera.
 
\paragraph{Dataset generation.} We built a separate single-class
dataset for each of the five categories. Keeping them separate at this
stage meant each dataset could be sourced, labelled, and checked on its
own. Every dataset uses the standard YOLO layout:
\texttt{train}/\texttt{val}/\texttt{test} splits, each with
\texttt{images} and \texttt{labels} directories and one normalized
\texttt{[class\_id, cx, cy, w, h]} label file per image. Before
training, we checked label quality with a visualization tool that draws
each label file's bounding boxes back onto its image, which made bad
annotations easy to spot. We then fine-tuned a YOLO11n detector on each
dataset.
 
\paragraph{Dataset merge.} An automated pipeline merges the five
single-class datasets into one five-class dataset. For every image it
keeps the original ground-truth labels, and it also runs every other
category's detector over that image to add cross-class pseudo-labels.
This matters because an image collected for one class often contains
unlabelled objects of another; without the pseudo-labels, the merged
model would be trained to treat those correct detections as mistakes.
Per-class confidence thresholds---0.30 for barrels, tires, and stop
signs, 0.35 otherwise---decide which pseudo-labels are kept. To keep
the merge practical on large datasets, it batches GPU inference,
overlaps file I/O with a thread pool, and writes an atomic checkpoint
after each batch so a crash does not mean starting over. The merged
\texttt{data.yaml} keeps the train/validation split intact. A single
YOLO11 model is then trained on the merged dataset, and those weights
are the detector we deploy on the vehicle.
 
\paragraph{ZED ROS 2 node.} On the vehicle, the merged model runs
inside the \texttt{yolo\_zed\_node} ROS~2 node. At startup the node
loads the model and reads its class set from a YAML file. Each frame,
it grabs a synchronized left RGB image and XYZ point cloud from the
ZED~2 and runs the image through the model. For each bounding box, it
samples a small point-cloud window around the box centre and takes the
median, minimum, and maximum range. NaN and non-finite returns are
dropped first, so a few bad pixels cannot throw off the reported
distance.
 
\paragraph{Internal representation.} The node publishes a
\texttt{YoloDetectionArray} message on \texttt{/yolo\_detections}. The
array has one fixed entry per model class, including classes not seen
in the current frame (their \texttt{detected} flag is false), so a
subscriber can index any class directly without checking whether the
field exists. Each entry holds the class name, confidence, pixel
bounding box, min/median/max range, and 3D position $(x, y, z)$ in
metres. The array is rebuilt from scratch every frame instead of being
accumulated, so it always shows what is currently visible---an object
drops out once it leaves the field of view and reappears on the first
frame it is detected again.
 
\paragraph{Performance.} \textit{we lwk did not take performance metrics, but will add them later if we have time}
 
\begin{table}[h]
\centering
\small
\begin{tabular}{p{3.6cm} p{2.4cm} p{2.4cm} p{4.0cm}}
\hline
\textbf{Requirement} & \textbf{Target} & \textbf{Measured} &
\textbf{What This Tells Us} \\
\hline
Merged model detection accuracy (mAP@0.5) &
[PLACEHOLDER] & [PLACEHOLDER] & [PLACEHOLDER] \\
\hline
Detection node frame rate &
[PLACEHOLDER] & [PLACEHOLDER] & [PLACEHOLDER] \\
\hline
Reported range vs.\ ground-truth distance &
[PLACEHOLDER] & [PLACEHOLDER] & [PLACEHOLDER] \\
\hline
\end{tabular}
\end{table}
